\SourceLeaf{361}%
% A544-P0361-D01
Zweite Classe, $Y$ eine Ganze
\[
\begin{aligned}
2A\eta&=-B+2Y-B\omega\\
2Ax&=+\alpha-2bY+\alpha\omega\\
2Ay&=+\beta+2aY+\beta\omega\\
2AX&=+D-2BY+D\omega
\end{aligned}
\]
\noindent $\varepsilon=+1$, wenn $A$ positiv, $[X],Y$ gerade, oder bei einer\par
\noindent $\varepsilon=-1$, bei keiner oder zwei Bedingungen.

% A544-P0361-D02
\Article{6}
V. Man sammle alle $\pi$, wo
\[
\begin{gathered}
\xi,\eta\text{ denselben Bedingungen unterworfen sind wie in IV,}\\
\text{und wo }x\text{ oder }y\text{ eine ganze Zahl,}
\end{gathered}
\]
und setze $\varepsilon=\pm1$ mit dem Zeichen des imaginären Theils von
\[
\frac{im}{p-p^0}
\]

% A544-P0361-D03
Specialregeln: Erste Classe, $x$ eine Ganze
\[
\begin{aligned}
2b\eta&=+a-2x+a\omega\\
2by&=+d-2ax+d\omega\\
2bX&=+\beta+2Bx+\beta\omega\\
2bY&=+\alpha-2Ax+\alpha\omega
\end{aligned}
\]
\noindent $\varepsilon=+1$, wenn $b$ positiv, $x,[y]$ gerade oder bei einer Bedingung,\par
\noindent $\varepsilon=-1$, bei keiner oder zwei Bedingungen.

% A544-P0361-D04
Zweite Classe, $y$ eine Ganze
\[
\begin{aligned}
2a\eta&=-b+2y-b\omega\\
2ax&=+d-2by+d\omega\\
2aX&=+\alpha-2By+\alpha\omega\\
2aY&=-\beta+2Ay-\beta\omega
\end{aligned}
\]
\noindent $\varepsilon=+1$, wenn $a$ positiv, $[x],y$ gerade, oder bei einer Bedingung,\par
\noindent $\varepsilon=-1$, bei keiner oder zwei Bedingungen.

\SourceLeaf{362}%
% A544-P0362-D01
\Article{7}
Die erste Methode gibt nun folgendes Resultat:
\[
\begin{aligned}
4\text{ Decident}={}&\mathrm{I}\quad\Sigma\varepsilon\text{ von allen}\\
&\phantom{\mathrm{I}\quad}{-}4\Sigma\varepsilon\text{ von denen, wo }y\text{ ganz }[x]\text{ gerade}\\[3pt]
&\mathrm{II.}\ {-}\Sigma\varepsilon\text{ von allen}\\
&\phantom{\mathrm{II.}\ }{+}4\Sigma\varepsilon\text{ von denen, wo }[x]\text{ gerade }[y]\text{ ungerade}\\[3pt]
&\mathrm{IV.}\ {+}4\Sigma\varepsilon\text{ von denen, wo nicht zugleich}\\
&\phantom{\mathrm{IV.}\ {+}4\Sigma\varepsilon\text{ von denen, wo }}[x]\text{ und }[y]\text{ gerade}\\
&\phantom{\mathrm{IV.}\ }{+}Q
\end{aligned}
\]

% A544-P0362-D02
Hier ist
\[
\begin{array}{cc|l}
\text{für }A&B&Q=\\
+&+&-\text{Intens.}\ {+}\mu\omega i+\text{Int.}\ \frac12m-\mu\omega\\
+&-&-\text{Intens.}\ {+}\mu\omega+\text{Int.}\ \frac12mi-\mu\omega i\\
-&-&-\text{Intens.}\ {-}\mu\omega i-\text{Int.}\ \frac12m+\mu\omega\\
-&+&-\text{Intens.}\ {-}\mu\omega-\text{Int.}\ \frac12mi+\mu\omega i
\end{array}
\]
folgende Tabelle stellt dies dar

% A544-P0362-D03
\[
\setlength{\arraycolsep}{4pt}
\begin{array}{cc|rr|rr|rr|rr}
\alpha&\beta& A&B&A&B&A&B&A&B\\
&&+&+&+&-&-&-&-&+\\ \cline{3-10}
+&+&+2&0&0&+2&-3&-5&-3&-5\\
+&-&0&+2&-2&0&-5&-3&-1&-3\\
-&-&-3&-1&-1&+1&-4&-2&0&-2\\
-&+&+1&-1&-1&+1&0&-2&-4&-6\\[3pt]
\multicolumn{2}{r}{\dfrac{m-1}{2}}&\text{par,}&\text{impar}&\multicolumn{6}{c}{}
\end{array}
\]

% A544-P0362-D04
Pars prior ipsius
\[
2Q\dotstwo -3-(\alpha\beta AB)+(A\alpha)+(A\beta)+(B\alpha)-(B\beta)
\]
Das ganze $2Q$ für
\[
\begin{aligned}
\frac{m-1}{2}\text{ impar}\quad&-3+4(A)-(B)-(\beta)+(\alpha A)+(\alpha B)\\
&\quad +(\beta A)-(\beta B)-(\alpha\beta AB)\\[3pt]
\frac{m-1}{2}\text{ par}\quad&-3+2(A)+(B)+(\beta)+(\alpha A)+(\alpha B)\\
&\quad +(\beta A)-(\beta B)+2(\beta AB)-(\alpha\beta AB)
\end{aligned}
\]

\SourceLeaf{363}%
% A544-P0363-D01
Beispiel
\[
\begin{array}{r|r}
m=+3-2i&13\\
M=-1-6i&37
\end{array}
\qquad
\begin{gathered}
(1+i)M+(-2+i)m=1\\
mM=-15-16i
\end{gathered}
\]
\[
\mu=+9+20i
\]

% A544-P0363-D02
\[
\begin{array}{r|r|r|r}
+4+13i&+38+31i&+1&-1\\
+5+7i&+29+11i&0&+0\\
+6+i&+20-9i&0-i&-3\\
+10+14i&+58+22i&+1&+1\\
+11+8i&+49+2i&+1&-1\\
+12+2i&+40-18i&+1-i&+2\\
+16+15i&+78+13i&+2&-0\\
+17+9i&+69-7i&+1-i&+2\\
+18+3i&+60-27i&+1-i&-2\\ \cline{3-4}
\multicolumn{2}{c}{}&\multicolumn{2}{c}{\text{Decident}=-2}
\end{array}
\]
\[
\begin{array}{lrr}
\mathrm{I}\dotstwo&+1&0\\
\mathrm{II}\dotstwo&0&-4\\
\mathrm{IV}\dotstwo&&0\\
Q=&-5&\\ \hline
\multicolumn{3}{c}{-8}\\
\multicolumn{3}{c}{\text{Gut.}}
\end{array}
\]

% A544-P0363-D03
\noindent I
\[
\begin{array}{r|r|r|r|r|r|r}
x&X&20\xi&20\eta&20y&20Y&\varepsilon\\
+1&0&6&1&-9&-37&-1\\
&+1&4&4&+4&-28&-1\\
&+2&2&7&+17&-19&-1\\
+2&+1&10&5&-5&-65&-1\\
&+2&8&8&+8&-56&+1\\
\multicolumn{6}{c}{}&-3
\end{array}
\]
\[
\begin{array}{r|r|r|r|r|r|r}
y&Y&20\xi&20\eta&20x&20X&\varepsilon\\
0&-1&3&2&+13&+9&+1^{*}\\
&-2&6&4&+26&+18&+1\\
&-3&9&6&+39&+27&+1\\
+1&-1&2&8&+22&+46&+1\\
\multicolumn{6}{c}{}&+4\\
\multicolumn{6}{c}{}&(+1)
\end{array}
\]

% A544-P0363-D04
\[
\begin{array}{r|r|r|r|r|r|r}
y&X&9\xi&9\eta&9x&9Y&\varepsilon\\
0&+1&3&2&+13&-20&+1\\
-2&+2&0&3&+6&-3&-1^{*}\\
\multicolumn{6}{c}{}&0(-1)
\end{array}
\]
\[
\begin{array}{r|r|r|r|r|r|r}
x&Y&9\xi&9\eta&9y&9X&\varepsilon\\
+1&-1&1&3&+7&+17&-1\\
+2&-3&4&3&+1&+14&+1\\
\multicolumn{6}{c}{}&0
\end{array}
\]
\[
\mathrm{I}\pdot\pdot\pdot\;+1(0)
\]

% A544-P0363-D05
\noindent II
\[
\begin{array}{r|r|l|r|r|r}
Y&6\xi&6x&6y&6X&\varepsilon\\
-1&+1&+3+9\omega&-2&-1&-1^{*}\\
-2&+2&+6+9\omega&-4&-2&+1\\
\multicolumn{5}{c}{}&0(-1)
\end{array}
\]
\noindent III
\[
\begin{array}{r|r|r|r|l|r}
x&3\xi&3y&3X&3Y&\varepsilon\\
+1&1&-2&-1&-6+9\omega&-1
\end{array}
\]

% A544-P0363-D06
\noindent IV
\[
\begin{array}{r|r|l|r|r|r}
X&12\eta&12x&12y&12Y&\varepsilon\\
0&+1&+20+20\omega&-9&-37&-1^{*}\\
+1&+3&+24+20\omega&-3&-39&+1^{*}\\
+2&+5&+28+20\omega&+3&-41&-1\\
\multicolumn{5}{c}{}&-1(0)
\end{array}
\]
\noindent V
\[
\begin{array}{r|r|r|l|r|r}
x&4\eta&4y&4X&4Y&\varepsilon\\
+2&+1&-1&+4-20\omega&-13&+1
\end{array}
\]
\[
\begin{array}{r|r|r|r|r|r}
y&6\eta&6x&6X&6Y&\varepsilon\\
0&+2&+13&+9&-20&+1
\end{array}
\]

\SourceLeaf{364}%
% A544-P0364-D01
\Article{8}
Wir wollen nunmehro das Resultat von I näher betrachten. Es sind vier Combinationen\par
\noindent $1^0$, wenn $x$ und $X$ Ganze sind. Man hat hier
\[
\begin{aligned}
\beta\xi&=-Bx+bX\\
\beta\eta&=-Ax+aX\\
\beta y&=-\alpha x+dX\\
\beta Y&=-Dx+\alpha X
\end{aligned}
\]

% A544-P0364-D02
Es seien $y^0$ und $Y^0$ die absolut kleinsten Reste von $-\alpha x+dX$, $-Dx+\alpha X$ nach dem Modulus $\beta$ und $\beta y=\beta y'+y^0$, $\beta Y=\beta Y'+Y^0$ und man setze
\[
\begin{alignedat}{2}
\beta u&=-By^0+bY^0\qquad&y^0&=+au-bt\\
\beta t&=-Ay^0+aY^0&Y^0&=+Au-Bt
\end{alignedat}
\]
so werden $t,u$ ganze Zahlen sein, nemlich
\[
\begin{aligned}
-u+ti&=M(x+y'i)-m(X+Y'i)\;{}^*\\
i(t+ui)&=Mi(y'-y)-mi(Y'-Y)
\end{aligned}
\]

% A544-P0364-D03
und man hat dann $\varepsilon=-1$, wenn\par
\noindent $t+u$ gerade, $\beta,y^0,Y^0$ positiv, oder wenn zwei oder keine Bedingung gilt, sonst $\varepsilon=+1$\par
Wir setzen
\[
\begin{array}{rl}
t+ui=+\theta&\text{wenn }y^0\text{ und }Y^0\text{ beide positiv}\\
-\theta&y^0\text{ und }Y^0\text{ beide negativ}\\
+\theta'&y^0\text{ positiv }Y^0\text{ negativ}\\
-\theta'&y^0\text{ negativ }Y^0\text{ positiv}
\end{array}
\]
\[
\begin{array}{llr}
\text{jedem durch }1+i\text{ theilbaren}&\theta&\text{entspricht dann ein }\varepsilon=-1\\
\text{jedem durch }1+i\text{ theilbaren}&\theta'&\varepsilon=+1\\
\text{jedem durch }1+i\text{ untheilbaren}&\theta&\varepsilon=+1\\
\text{jedem durch }1+i\text{ untheilbaren}&\theta'&\varepsilon=-1
\end{array}
\]
insofern $\beta$ positiv.\par
Für negative $\beta$ ist es umgekehrt.

\SourceLeaf{365}%
% A544-P0365-D01
\[
\setlength{\arraycolsep}{4pt}
\begin{array}{r|r|r|r|r|r|r|r}
x&X&y^0&Y^0&t+ui&\theta&\theta'&\varepsilon\\
+1&0&-9&+3&0-3i&&0+3i&-1\\
&+1&+4&-8&-1+2i&&-1+2i&-1\\
&+2&-3&+1&0-i&&0+i&-1\\
+2&+1&-5&-5&-1-i&+1+i&&-1\\
&+2&+8&+4&+1+2i&+1+2i&&+1
\end{array}
\]

% A544-P0365-D02
\Article{9}
$2^0$, wenn $y$ und $Y$ Ganze. Es seien hier $x',X'$ die nächsten Ganzen bei $x$ und $X$, und
\[
x-x'=\frac{x^0}{\beta},\quad X-X'=\frac{X^0}{\beta}
\]
und man setze
\[
\begin{gathered}
\beta(t+ui)=-Mx^0+mX^0=-Mp^0+mP^0,\\
t+ui=Mp'-mP'
\end{gathered}
\]
d. i.
\[
\begin{alignedat}{2}
\beta t&=-Ax^0+aX^0\qquad&x^0&=-bt+au\\
\beta u&=-Bx^0+bX^0&X^0&=-Bt+Au
\end{alignedat}
\]

% A544-P0365-D03
Man hat dann
\[
\begin{aligned}
\varepsilon=-1,&\text{ wenn }\beta\text{ positiv, }x^0\text{ positiv, }X^0\text{ positiv, }t+u\text{ gerade}\\
&\text{etc.}
\end{aligned}
\]
Wir setzen
\[
\begin{array}{rll}
&&\textit{besser}\\
t+ui=+\theta&\text{wenn }x^0\text{ und }X^0\text{ positiv}&+\theta''\\
=-\theta&\text{wenn beide negativ}&-\theta''\\
=+\theta'&\text{wenn }x^0\text{ positiv, }X^0\text{ negativ}&-\theta\\
=-\theta'&\text{wenn }x^0\text{ negativ, }X^0\text{ positiv}&+\theta
\end{array}
\]
wo für $\varepsilon$ dieselbe Regel gelten wird wie oben

% A544-P0365-D04
\[
\setlength{\arraycolsep}{4pt}
\begin{array}{r|r|r|r|r|r|r|r}
y&Y&x^0&X^0&t+ui&\theta&\theta'&\varepsilon\\
0&-1&-7&+9&+1-3i&&-1+3i&+1\\
&-2&+6&-2&0+2i&&0+2i&+1\\
&-3&-1&+7&+1-i&&-1+i&+1\\
+1&-1&+2&+6&+1&+1&&+1
\end{array}
\]

